%-----------------------------------------------
%  Engineer's & Master's Thesis Template
%  Copyleft by Artur M. Brodzki & Piotr Woźniak
%  Warsaw University of Technology, 2019-2022
%-----------------------------------------------

\documentclass[
    bindingoffset=5mm,  % Binding offset
    footnoteindent=3mm, % Footnote indent
    hyphenation=true    % Hyphenation turn on/off
]{src/wut-thesis}

\graphicspath{{tex/img/}} % Katalog z obrazkami.
\addbibresource{bibliografia.bib} % Plik .bib z bibliografią

\facultyeiti    % Wydział Elektroniki i Technik Informacyjnych
\MasterThesis % Praca inżynierska
\langpol % Praca w języku polskim

\begin{document}

%------------------
% Strona tytułowa
%------------------
\instytut{Instytut telekomunikacji}
\kierunek{Telekomunikacja}
\specjalnosc{Teleinformatyka i Zarządzanie w Telekomunikacji}
\title{
    Sworzenie platformy do modelowania i uruchamiania zamkniętych pętli sterowania w Kubernetes
}
\engtitle{
    Creation of a platform for designing and running closed control loops in Kubernetes
}

\author{Andrzej Gawor}
\album{300528}
\promotor{dr inż. Dariusz Bursztynowski}
\date{\the\year}
\maketitle

%-------------------------------------
% Streszczenie po polsku 
%-------------------------------------
\cleardoublepage % Zaczynamy od nieparzystej strony
\abstract
Projekt opisany w niniejszej pracy skupia się na zaproponowaniu oraz przedstawieniu implementacji  architektury platformy, która pozwala na modelowanie oraz uruchamianie zamkniętych pętli sterowania w Kubernetes. Genezą projektu jest praca jednego z komitetów ETSI o nazwie "ENI - Experiential Networked Intelligence", która skupia się na ułatwieniu pracy operatora sieci telekomunikacyjnych wykorzystując mechanizmy sztucznej inteligencji w zamkniętych pętlach sterowania. ENI w jednym ze swoich dokumentów dokonuje przeglądu zamkniętych pętli sterowania znanych ludzkości z innych dziedzin. 

Naturalnym następnym krokiem jest zapropowanie platformy, na której operator mógłby takowe pętle projektować oraz uruchamiać. W tym celu zdefiniowanio zestaw wymagań oraz założeń dla takiego systemu. Jako środowisko uruchomieniowe wybrano  Kubernetes z racji, że jest to system dobrze znany w społeczności oraz sam natywnie używa zamkniętych pętli sterowania. Następnie przeprowadzono obszerną analizę jak za pomocą mechanizmów rozszerzania Kubernetes takich jak "Custom Resources" oraz "Operator" pattern można stworzyć framework umożliwiający modelowanie zamkniętych pętli sterowania. Praca opisuje powstałą platformę, jej architekturę, semantykę składni w definiowanych obiektach, zasady działania, integracje z zewnętrznymi systemami oraz instrukcję jej użytkowania. Omówiona została również implemtacja platformy, technologie za nią stojące oraz decyzje podjęte podczas jej powstawania. Finalnie przedstawiono również test działania platformy w praktyce wykorzystując do tego emulator systemu 5G jakim jest Open5GS w połączeniu z UERANSIM. Pracę podsumuję lista wniosków oraz potencjalnych dróg rozwoju platformy. 
\keywords Zamknięte pętle sterowania, Kubernetes, Zarządzanie sieciami telekomunikacyjnymi, Automatyzacja, Go, Open5GS, Mikroserwisy

%----------------------------------------
% Streszczenie po angielsku
%----------------------------------------
\clearpage
\secondabstract
The project described in this thesis focuses on proposing and implementing a reusable architecture (hereafter referred to as "the Platform") that enables the modeling and execution of closed control loops in Kubernetes. The genesis of the project lies in the work of one of the ETSI committees called "ENI - Experiential Networked Intelligence," which aims to simplify the work of telecommunications network operators by leveraging artificial intelligence mechanisms in closed control loops based on metadata-driven and context-aware policies. In one of its documents, ENI reviews closed control loops known to humanity from other fields. A natural next step is to propose a platform on which operators could design and execute such loops. 

To achieve this, a set of requirements and assumptions for such a system was defined. Kubernetes was chosen as the runtime environment due to its widespread adoption in the community and its inherent use of closed control loops. An extensive analysis was conducted on how Kubernetes extension mechanisms, such as "Custom Resources" and the "Operator" pattern, could be used to create a framework enabling the modeling of closed control loops. 

This thesis describes the developed platform, its architecture, the semantics of syntax in the defined objects, operational principles, integrations with external systems, and a user guide. It also discusses the platform's implementation, the technologies behind it, and the decisions made during its development. Finally, the thesis presents a practical test of the platform's functionality using the Open5GS 5G system emulator in combination with UERANSIM. The work concludes with a list of findings and potential extensions or improvements to the platform.
\secondkeywords Closed Control Loops, Kubernetes, Managing Telco-Networks, Automation, Go, Open5GS, Microservices

\pagestyle{plain}

%--------------
% Spis treści
%--------------
\cleardoublepage % Zaczynamy od nieparzystej strony
\tableofcontents

%------------
% Rozdziały
%------------
\cleardoublepage % Zaczynamy od nieparzystej strony
\pagestyle{headings}

\newpage
\section{Wstęp}
\subsection{Przedmowa}


\subsection{Cel i zakres pracy}

Celem pracy jest kontynuacja badań ENI, polegająca na zaproponowaniu \hyperlink{def:architektura}{\textit{architektury}} platformy, na której możliwe będzie modelowanie oraz uruchamianie zamkniętych pętli sterowania. Niniejsza praca nie skupia się na aspektach związanych ze sztuczną inteligencją. Platforma ma jedynie służyć do zamodelowania oraz uruchomienia \hyperlink{def:workflow}{\textit{workflow}} pętli, ale sama nie stanowi środowiska wykonawczego dla jej komponentów. 

W zakres pracy wchodzi: sformułowanie wymagań na platformę, opracowanie proponowanej architektury wsparte licznymi badaniami podczas jej powstawania, implementacja PoC (ang. \textit{Proof of Concept}) według proponowanej architektury, przeprowadzenie testów platformy oraz analiza jej potencjału w kontekście dalszego rozwoju. Po publikacji praca może stanowić podstawę do implementacji . \hyperlink{def:kognitywny-system-zarzadzania-siecia}{Kognitywnych Systemów Zarządzania Siecią}, zgodnych ze specyfikacjami ENI. 
\subsection{Struktura pracy}

Praca została podzielona na 6 rozdziałów. Drugi rozdział przedstawia bardziej szczegółowo niż we wstępie badania podjęte przez ENI. Stanowi on wprowadzenie teoretyczne oraz pojęciowe, aby ułatwić przekaz w dalszej części pracy. Rozdział trzeci zawiera opis proponowanej architektury. Rozdział czwarty opisuje implementację PoC platformy oraz napotkane wyzwania podczas formułowania architektury, które celowo zostały zebrane w jedno miejsce i umieszczone oddzielnie w celu łatwiejszej lektury. Rozdział numer pięć przedstawia użycie platformy w praktyce przy okazji stanowiąc jej test. Na koniec, w rozdziale szóstym przeprowadzono analizę zaproponowanej architektury, jej możliwości oraz ograniczenia oraz wskazano kierunki potencjalnego rozwoju. 
%---------------
% Bibliografia
%---------------
\cleardoublepage % Zaczynamy od nieparzystej strony
\printbibliography
\clearpage

% Wykaz symboli i skrótów.
% Pamiętaj, żeby posortować symbole alfabetycznie
% we własnym zakresie. Makro \acronymlist
% generuje właściwy tytuł sekcji, w zależności od języka.
% Makro \acronym dodaje skrót/symbol do listy,
% zapewniając podstawowe formatowanie.
\acronymlist
\acronym{EiTI}{Wydział Elektroniki i Technik Informacyjnych}
\acronym{PW}{Politechnika Warszawska}
\acronym{WEIRD}{ang. \emph{Western, Educated, Industrialized, Rich and Democratic}}
\vspace{0.8cm}

%--------------------------------------
% Spisy: rysunków, tabel, załączników
%--------------------------------------
\pagestyle{plain}

\listoffigurestoc    % Spis rysunków.
\vspace{1cm}         % vertical space
\listoftablestoc     % Spis tabel.
\vspace{1cm}         % vertical space
\listofappendicestoc % Spis załączników

%-------------
% Załączniki
%-------------

% Obrazki i tabele w załącznikach nie trafiają do spisów
\captionsetup[figure]{list=no}
\captionsetup[table]{list=no}

% Załącznik 1
\clearpage
\appendix{Nazwa załącznika 1}
\lipsum[1-3]
\begin{figure}[!h]
	\centering \includegraphics[width=0.5\linewidth]{logopw2.png}
	\caption{Obrazek w załączniku.}
\end{figure}
\lipsum[4-7]

% Załącznik 2
\clearpage
\appendix{Nazwa załącznika 2}
\lipsum[1-2]
\begin{table}[!h] \centering
    \caption{Tabela w załączniku.}
    \begin{tabular} {| c | c | r |} \hline
        Kolumna 1       & Kolumna 2 & Liczba \\ \hline\hline
        cell1           & cell2     & 60     \\ \hline
        \multicolumn{2}{|r|}{Suma:} & 123,45 \\ \hline
    \end{tabular}
\end{table}
\lipsum[3-4]

% Używając powyższych spisów jako szablonu,
% możesz dodać również swój własny wykaz,
% np. spis algorytmów.

\end{document} % Dobranoc.
